% !TEX root = /home/eb/.vscode/extensions/qftrules.programme-colle/tmp/Exercice.tex

%%%%%%%%%%%%%%%%%%%%%%%%%%%%%%%%%%%%%%%%%%%%%%%%%%%%%%%%%%%%%%%%
\begin{exo}[PC][3][python]{Mesure d'une distance focale et simulation Monte-Carlo}
\Notions{\item Capacité numérique \task Incertitudes}

Voir la \refCN{4} sur capytale :
\begin{center}
	\href{https://capytale2.ac-paris.fr/web/c/c29d-2482234}{https://capytale2.ac-paris.fr/web/c/c29d-2482234}
\end{center}

% En montage d'optique, on utilise souvent des lentilles comme système de projection. Pour interpréter l'expérience, il faut donc connaître précisément la distance focale $f'$ de la lentille. On peut pour cela mesurer les positions d'un objet $A$ sur l'axe optique et de son image $A'$ puis appliquer la formule de Descartes
% \begin{equation}\label{eq:descartes}
% \frac{1}{f'} = \frac{1}{\overline{OA'}} - \frac{1}{\overline{OA}}.
% \end{equation}
% Seulement, il n'y a pas de formule simple pour déterminer l'incertitude sur $f'$ pour la fonction compliqué de l'\Eq{descartes}. On recourt alors à une simulation de type \imp{Monte-Carlo} pour étudier la variabilité de la grandeur $f'$ sans avoir à composer d'incertitudes.
% \begin{quest}
% \item Définir en \texttt{python} une fonction \texttt{focale} qui retourne $f'$ connaissant $\overline{OA}$ et $\overline{OA'}$.


% \sol{
% \python{monte_carlo_0_sol}
% }
% \item On associe maintenant une distribution de probabilité aux grandeurs mesurées $\overline{OA}$ et $\overline{OA'}$. On suppose cette distribution uniforme, centrée sur  $\SI{-15}{cm}$ et $\SI{30}{cm}$ respectivement, et de demi-largeur égale à $\SI{1}{cm}$, selon le code suivant.
% \python[linenos=true]{monte_carlo_1}
% Compléter le code ci-dessus pour réaliser $N=10000$ tirages aléatoires des grandeurs $\overline{OA}$ et  $\overline{OA'}$, et de stocker les valeurs de $f'$ correspondantes sous forme de liste. On pourra pour cela utiliser la fonction \texttt{np.random.uniform}.

% \sol{
% \python{monte_carlo_1_sol}
% }

% \item Utiliser la fonction \texttt{hist} de la bibliothèque \texttt{matplotlib.pyplot} pour afficher un histogramme des valeurs de $f'$ comprenant $40$ boîtes (option \texttt{bins} de la fonction \texttt{hist}). La figure comportera un titre et un libellé pour les abscisses.

% \sol{
% \python{monte_carlo_2_sol}
% \includegraphics[scale=1]{distrib_f.pdf}
% }

% \item Enfin, afficher la valeur moyenne et l'écart-type de $f'$. Répéter plusieurs fois la simulation et comparer les valeurs de $f'$ obtenues. Comparer au cas avec $N=10$ tirages.
% \sol{
% \python{monte_carlo_3_sol}

% On trouve $f'=\SI{10.0\pm0.3}{cm}$. La valeur de $f'$ et de $\Delta f'$ varie à $\num{0.001}$ près pour $N=10000$. Au contraire, pour $N=10$, les fluctuations sont beaucoup plus grandes, de l'ordre de $\num{0.1}$. Cela montre l'intérêt de faire un maximum de points de mesure pour avoir un résultat de mesure fiable.
% }
% \end{quest}
\end{exo}
%%%%%%%%%%%%%%%%%%%%%%%%%%%%%%%%%%%%

%%%%%%%%%%%%%%%%%%%%%%%%%%%%%%%%%%%%%%%%%%%%%%%%%%%%%%%%%%%%
\begin{exo}[MPSI][2][colle]{Points de Weierstrass}
On considère un dioptre sphérique de centre C, de sommet S, d'axe optique (CS), et A un point à l'intérieur de la sphère. Le milieu à l'intérieur de la sphère est d'indice de réfraction~$n$, le milieu extérieur d'indice de réfraction $n'$. On souhaite déterminer à quelle condition le dioptre sphérique est rigoureusement stigmatique pour le point A.

%	%%% FIGURE %%%
%	\begin{figure}[h!]
%	\centering
\begin{center}
	\includegraphics[scale=1]{dioptre.pdf}
\end{center}
%	\caption{Dioptre sphérique. On note S un des sommets du dioptre, et C son centre.}
%	\label{fig:}
%	\end{figure}
%	%%%%%%%%%%

\begin{quest}

	\item Le point image A' est-il réel ou virtuel ? Sur le schéma de la Figure 2, lequel des deux indices $n$ ou $n'$ est le plus grand ?

	\sol{virtuel, $n>n'$}
\end{quest}
On admet qu'un système optique est rigoureusement stigmatique pour un couple de points A et A' si le temps $\tau_{\rm AA'}$ de parcours de la lumière entre ces deux points ne dépend pas du trajet du rayon. Pour un trajet virtuel, le temps de parcours est compté négativement et l'indice de réfraction vu par le rayon est celui du milieu en sortie. Soit I le point de contact du rayon sur le dioptre et H son projeté orthogonal sur l'axe optique (CS).
\begin{quest}
	\item  Démontrer la relation $AI=yA'I$ avec $y=n'/n$ le rapport des indices.

	\sol{$AI/A'I=n'/n$}

	\item On note $R=\overline{CS}$ le rayon du dioptre, $x=\overline{HS}$, $p=\overline{AC}$ et $p'=\overline{A'C}$. En utilisant le théorème de Pythagore, exprimer $\overline{AI}^2$ et $\overline{A'I}^2$ en fonction de $p$, $R$ et $x$. 
	% On pourra penser à décomposer les distances algébriques de sorte à faire apparaître $x$, $R$ et $p$. 
	% Faire de même pour $\overline{A'I}^2$.

	\sol{$\overline{AI}^2=(R+p)^2+2xp$}

	\item Sachant que la relation $AI=yA'I$ ne dépend pas de la position du point $\rm I$, trouver deux relations entre $p$ et $p'$. En déduire l'expression de $p$ et $p'$ en fonction de $R$ et $y$, et donc la position des points A et~A'. Ces points conjugués sont appelés points de Weierstrass.

	\item Dans le cas particulier ou la droite (CI) est orthogonale à (CS), vérifier que la loi de Snell-Descartes de la réfraction est vérifiée.

\end{quest}
\end{exo}
%%%%%%%%%%%%%%%%%%%%%%%%%%%%%%%%%%%%%%%%%%%%%%%%%%%%%%%%%%%%


%%%%%%%%%%%%%%%%%%%%%%%%%%%%%%%%%%%
\begin{exo}[PC][1][oral]{Détermination d’une distance focale}
	On souhaite étudier la photographie ci-dessous.
La mise au point est effectuée sur la figurine qui se situe à une distance de $40 \mathrm{~cm}$ de l'objectif.

\begin{fig}[label]
	\centering
	\includegraphics[width=0.7\columnwidth]{playmobil.jpg}
	%\capt{}
\end{fig}

Les réglages effectués sur l'appareil pour prendre cette photographie sont les suivants :
Vitesse d'obturation : $1/500$; Diamètre d'ouverture : $f/\num{4.0}$ ; ISO : $160$, où
$f$ correspond à la valeur de la distance focale de l'objectif. Cette photographie a été prise avec un appareil photographique dont les caractéristiques générales sont les suivantes.
\begin{center}
	
	\begin{tabular}{|l|l|}
		\hline Pixels effectifs de l'appariel photo (Max.) & Environ 20,3 millions de pixels \\
		\hline Distance focale de l'objectif & de 25 à $750 \mathrm{~mm}$ \\
		\hline Dimensions du capteur photosensible & $24 \times 36 \mathrm{~mm}$ \\
		\hline
	\end{tabular}
\end{center}

À l'aide des informations fournies, estimer la valeur de la distance focale de l'objectif utilisé pour prendre la photographie. En déduire la valeur du diamètre d'ouverture correspondant.

\sol{
On demande de déterminer une distance focale $f$. On se place dans le cadre de l’\udl{optique géométrique}, dans les \udl{conditions de Gauss}. Les autres distances données dans le problème sont : la distance $d=\SI{40}{mm}$ entre l’objectif et l’objet, et les dimensions $24\times36\SI{}{mm}$ de la cellule photosensible, le diamètre d’ouverture $D=f/4$. \\

La hauteur réelle de l’objet  peut s’estimer à $AB\simeq \SI{5}{cm}$ (taille typique d’un playmobile).
La hauteur $A’B’$ de l’image par l’objectif peut s’estimer à la moitié de la hauteur de la cellule photosensible en inspectant la photographie, soit $A’B’ \simeq \SI{12}{mm}$. \\

Le schéma optique qui modélise la situation est le suivant.

Le grandissement a pour expression
$$
\gamma  = \frac{\overline{OA’}}{\overline{OA}}= \frac{\overline{A’B’}}{\overline{AB}} = -\num{0.24}.
$$
Attention au signe ! La relation de conjugaison de Descartes donne ensuite,
$$
\frac{1}{\overline{OA’}}-\frac{1}{\overline{OA}}  = \frac{1}{\overline{OA’}}(1-\gamma) = \frac{1}{f} \Rar f = \frac{\overline{OA’}}{1-\gamma},
$$
soit 
\eqbox{
f = \frac{d}{1-\gamma} \simeq \SI{32}{cm}.
}
}
\end{exo}
%%%%%%%%%%%%%%%%%%%%%%%%%%%%%%%%%%%

%%%%%%%%%%%%%%%%%%%%%%%%%%%%%%%%%%%%%%%%%%%%%%%%%%%%%%%%%%%%
\begin{exo}[MPSI]{Les points de Weierstrass}
On considère un dioptre sphérique de centre C, de sommet S, d'axe optique (CS), et A un point à l'intérieur de la sphère. Le milieu à l'intérieur de la sphère est du verre d'indice de réfraction~$n$, le milieu extérieur est de l’air d'indice de réfraction $n_{\rm air}=1$. On souhaite déterminer à quelle condition le dioptre sphérique est rigoureusement stigmatique pour le point A.

%	%%% FIGURE %%%
%	\begin{figure}[h!]
%	\centering
\begin{center}
	\includegraphics[scale=1]{weier.pdf}
\end{center}
%	\caption{Dioptre sphérique. On note S un des sommets du dioptre, et C son centre.}
%	\label{fig:}
%	\end{figure}
%	%%%%%%%%%%

\begin{quest}

	\item Que signifie « stigmatisme rigoureux » ? Le point image A' est-il réel ou virtuel ? Sur le schéma ci-dessus, justifier que le rayon réfracté en I s’éloigne de la normale.

	\sol{virtuel, $n>n_{\rm air}$}

	\item On admet qu'un système dioptrique est rigoureusement stigmatique pour un couple de points A et A' si le temps de parcours $\tau_{\rm AA'}$  de la lumière entre ces deux points est nul quel que soit le trajet du rayon. Pour un trajet virtuel, le temps de parcours est compté négativement, et l'indice de réfraction vu par le rayon est celui du milieu en sortie. Montrer que $n AI=A'I$ pour tout I.

	\sol{$\tau_{AA’} = (n/c)AI - (1/c)AI’ = \cte$ donc $AI/A'I=1/n$.}


	\item On cherche la relation de conjugaison du dioptre sphérique. Soit H le projeté orthogonal de I sur l'axe optique (CS). Soit $\gamma$ l’angle $(ACI)=(A’CI)$. Exprimer $\sin(\gamma) $ de deux façons, puis en utilisant la condition de stigmatisme et la loi de Snell-Descartes, montrer que $A’C = n^2 AC$.


	\sol{
	On a $n \sin(i) = \sin(r)$ et la loi des sinus donne $\sin(\gamma) = (AI/p)\sin(i) = (A’I)/p’ \sin(r) = A’I/(np) \sin(i) = A’I/(n^2 p) \sin(r)$ donc $p’ = n^2 p$.
	}
%
%	\item Sachant que la relation $nAI=A'I$ ne dépend pas de la position de I (et donc de $x$), trouver deux relations entre $p$ et $p'$. En déduire l'expression de $p$ et $p'$ en fonction de $R$ et $n$, et donc la position des points A et~A'. Ces points conjugués sont appelés points de Weierstrass.

%	\item Dans le cas particulier ou la droite (CI) est orthogonale à (CS), vérifier que la loi de Snell-Descartes de la réfraction est vérifiée.

\end{quest}
\end{exo}
%%%%%%%%%%%%%%%%%%%%%%%%%%%%%%%%%%%%%%%%%%%%%%%%%%%%%%%%%%%%

%%%%%%%%%%%%%%%%%%%%%%%%%%%%%%%%%%%%%%%%%%%%%%%%%%%%%%%%%%%%%
\begin{exo}[PCSI]{Image par une lentille convergente - 1}
Un lentille mince convergente, de centre $O$ et de distance focale $f’=\SI{20}{cm}$, fait l’image d’un objet transversal AB de taille $\overline{AB}=\SI{5}{cm}$ situé à une distance $\overline{OA}=-\SI{25}{cm}$.
\begin{quest}
\item Faire un schéma de la situation et tracer les rayons lumineux. L’objet est-il réel ou virtuel ? Même question pour l’image.
\item Déterminer la distance $\overline{F’A’}$ de deux façons différentes (on utilisera les relations de conjugaison).
\item Définir le grandissement transversal $\gamma$. Le calculer.
\end{quest}

\end{exo}
%%%%%%%%%%%%%%%%%%%%%%%

%%%%%%%%%%%%%%%%%%%%%%%%%%%%%%%%%%%%%%%%%%%%%%%%%%%%%%%%%%%%
\begin{exo}[MPSI][1][colle]{Théorie de Descartes de l'arc-en-ciel}
Une goutte d'eau sphérique d'indice de réfraction $n\simeq\SI{1.33}{}$ est éclairée par le Soleil. On suppose que l'air environnant est d'indice $n_{\rm air}\simeq1$. Soit $D_N$ la déviation angulaire totale d'un rayon pénétrant la goutte avec un angle d'incidence $i$ et ressortant de la goutte après $N$ réflexions internes partielles. $D_N$ est donc la différence entre l'angle d'entrée et l'angle de sortie.


%%%% FIGURE %%%
%\begin{figure}[h!]
%\centering
\begin{center}
\includegraphics[scale=1]{goutte.pdf}
\end{center}
%\caption{Trajet d'un rayon lumineux frappant une goutte d'eau sphérique (seules une réflexion interne est représentée). Le rayon incident frappe la goutte au point A de la droite vers la gauche, à l'horizontale. Les traits pointillés représentent la normale au dioptre formé par la surface de la goutte.}
%\label{fig:goutte}
%\end{figure}
%%%%%%%%%%


\begin{quest}
	\item À l'aide des lois de Snell-Descartes et de considérations géométriques simples, justifier que les angles d'incidence $i$ et de réfraction $r$ définis au point de rencontre du rayon incident avec la goutte, noté A, se retrouvent en B et C comme indiqué sur la figure. On ne demande pas de déterminer $r$ à ce stade.

	\sol{Cercle donc triangles isocèles et loi de la réflexion.}
	\item Montrer qu'un rayon qui subit une réflexion interne subit une déviation angulaire $D_1=\pi-2r+2(i-r)$.   Généraliser au cas d'un rayon qui subit $N$ réflexions internes.

	\sol{Réflexion interne : $\pi-2r$; réfraction : $i-r$.

	%, et exprimer $D_N$ en fonction de $i$, $r$ et $N$.
$D_N=2(i-r)+N(\pi-2r)$.
}

%	\item En utilisant les lois de Snell-Descartes, exprimer l'angle de réfraction $r$ en fonction de l'angle d'incidence $i$. Montrer l'existence d'un extremum de $D_N$ (noté $D_N^m$) pour une certaine valeur $i_N^m$ de l'angle d'incidence. On donnera l'expression de $\cos^2(i_N^m)$. On rappelle qu'une fonction $x \mapsto f(x)$ est extrémale lorsque $f'(x)=0$.
%
%	\sol{$\cos^2(i_N^m) = (n^2-1)/(N^2+2N)$}

	\item En utilisant les lois de Snell-Descartes, exprimer l'angle de réfraction $r$ en fonction de l'angle d'incidence $i$. Exprimer la dérivée de $r$ par rapport à $i$, notée $\dd{r}{i}$. Montrer ensuite que
	% $$
	\begin{equation}
		\dd{D_1}{i} = 2-\frac{4}{n} \frac{\cos i}{\cos r}.
	\end{equation}
	% $$

	\sol{
$\dd{r}{i} = 	 \frac{1}{n} \frac{\cos i}{\cos r}$.
	}

%	Montrer l'existence d'un extremum de $D_N$ pour une certaine valeur $i_2^m$ de l'angle d'incidence. On donnera l'expression de $\cos^2(i_N^m)$. On rappelle qu'une fonction $x \mapsto f(x)$ est extrémale lorsque $f'(x)=0$.

	\item Montrer que $\dd{D_1}{i}$ s’annule lorsque
	% $$
	% $$
	\begin{equation}
		\sin^2(i) = \frac{4-n^2}{3}.
	\end{equation}
	En déduire que les rayons lumineux se concentrent autour d’une direction bien précise par rapport à la position du Soleil. Comment expliquer que les couleurs soient séparées dans un arc-en-ciel ?
%
%	\item Interpréter l'existence de cet extremum du point de vue de la quantité de lumière reçue par un observateur au sol. Calculer les valeurs numériques de $i_N^m$ et $D_N^m$ pour $n=1.33$, avec $N=1$ puis $N=2$.
%
%	\sol{$N=1$: $i_1^m=59 \degres$ et $D_1^m=138 \degres$, $N=2$, $i_1^m=72 \degres$ et $D_1^m=231 \degres$. L'extremum correspond à une accumulation de lumière.}
%
%	\item Comment expliquer que les couleurs soient séparées dans un arc-en-ciel ?
%
	\sol{dispersion}

\end{quest}
\end{exo}
%%%%%%%%%%%%%%%%%%%%%%%%%%%%%%%%%%%%%%%%%%%%%%%%%%%%%%%%%%%%

%%%%%%%%%%%%%%%%%%%%%%%%%%%%%%%%%%%%%%%%%%%%%%%%%%%%%%%%%%%%
\begin{exo}[MPSI]{Relation de Gullstrand}
Le but de l'exercice est de déterminer la vergence équivalente $V$ d'un système centré constitué de deux lentilles minces $\mc{L}_1$ et $\mc{L}_2$. Les focales images des lentilles sont $f'_1$ et $f'_2$, les vergences $V_1$ et $V_2$, et la distance $\overline{O_1 O_2} >0$ séparant les centres des lentilles est notée $e$.
%%%% FIGURE %%%
%\begin{figure}[h!]
%\centering
\begin{center}
\includegraphics[scale=1]{gull.pdf}
\end{center}
%\caption{Association de deux lentilles.}
%\label{fig:}
%\end{figure}
%%%%%%%%%%%
Pour définir la distance focale équivalente d'un système optique, il faut au préalable définir le point principal image $H'$, de la façon suivante. Un rayon parallèle à l'axe optique arrive sur $L_1$, traverse $L_2$ puis coupe l'axe optique en $A'$. On note $G$ l'intersection de la prolongation des rayons incident et émergent. Le point $H'$ est alors dans le plan orthogonal à l'axe optique passant par $G$ et situé sur l'axe optique. 
\begin{quest}
	\item Tracer la marche de ce rayon lumineux ainsi que les points $H',A'$ et $G$ dans le cas où les deux lentilles sont convergentes.

	\item Écrire la relation de conjugaison pour la lentille $\mc{L}_2$ faisant intervenir $A'$ avec origine au centre, en introduisant la vergence $V_2$.
	% On introduira la vergence $V_2=1/f'_2$.

	\sol{$1/\overline{O_2 A'} - 1/\overline{O_2 F_1'} = 1/f'_2 = V_2$.}

	\item Soit $D$ la distance séparant l'axe optique du point d'impact du rayon incident sur $L_1$, noté $I$. Soit $d$ la distance séparant l'axe optique du point d'impact du rayon sur $\mc{L}_2$, noté $J$, et $\alpha$ l'angle $(\overrightarrow{F'_1 0_1}, \overrightarrow{F'_1 I})$. Dans les conditions de Gauss, exprimer $\alpha$ de deux façons différentes pour trouver $d$.

	 \sol{ $\tan \alpha \approx \alpha \approx d/\overline{O_2 F_1'} \approx DV_1$ d'où $d= \overline{O_2 F_1'} DV_1$.}

	\item La distance focale du système optique est définie par $f' = 1/V = \overline{H' A'}$. Exprimer $\overline{O_2 A'}$ en fonction de la vergence~$V$.

	\sol{Par thalès, $\overline{O_2 A'}=d/(DV)=V_1 \overline{O_2 F'_1} / V$.}

	\item Combiner les deux équations trouvées en 3) et 4) pour exprimer $\overline{O_2 F'_1}$ puis déterminer $V$ en fonction de $V_1$, $V_2$ et $e$.  Cette relation est appelée relation de Gullstrand. Commenter physiquement cette relation (influence de la distance entre les lentilles, de leur caractère convergent ou divergent, etc.)

	\sol{$\overline{O_2 F'_1} = 1/V_1 - e$ d'où $V=V_1 + V_2 -e V_1 V_2$.}

\end{quest}
\end{exo}
%%%%%%%%%%%%%%%%%%%%%%%%%%%%%%%%%%%%%%%%%%%%%%%%%%%%%%%%%

%%%%%%%%%%%%%%%%%%%%%%%%%%%%%%%%%%%%%%%%%%%%%%%%%%%%%
\begin{exo}[MPSI][2][colle]{Spectroscope à prisme}
Un spectroscope à prisme est un appareil destiné à étudier le spectre d'une source lumineuse. 
% Un collimateur permet de réaliser un faisceau de rayons parallèles qui va éclairer un prisme. Un viseur permet ensuite d'étudier la lumière ayant traversé le prime.
Le prisme utilisé est d’indice optique $n$ et d’angle au sommet $A$. Sa section droite est un triangle d'angle au sommet $A$. L’indice de l’air est pris égal à $1$. Un rayon incident rencontre la face d'entrée au point $I$ sous l'angle d'incidence $i$, et émerge par l'autre face au point $J$ sous l'angle $i'$. On définit aussi les angles $r$ et $r’$ sur le schéma ci-dessous.

% On suppose jusqu'à la question 4 la lumière monochromatique et l'indice du prisme sera égal à $n$.


%%%% FIGURE %%%
%\begin{figure}[h!]
%\centering
\begin{center}
\includegraphics[scale=1.2]{prisme.png}
\end{center}
%\caption{Trajet d'un rayon lumineux dans un prisme (le point noté $I'$ sur le schéma correspond au point $J$ de l'énoncé).}
%\label{fig:}
%\end{figure}
%%%%%%%%%%%


\begin{quest}
\item En appliquant les lois de Descartes, écrire la relation entre $i$ et $r$ d'une part, et $i'$ et $r'$ d'autre part. On suppose que le prisme permet d'existence du rayon émergeant et on néglige, dans la suite, tout réflexion. Trouver une relation entre $r$, $r'$ et $A$.

\sol{Les démonstrations sont celles du cours, il faut écrire :
\begin{equation*}
\sin i = n \sin r \qquad n \sin r' = \sin i' \qquad A = r + r'
\end{equation*}}

\item Soit $D$ l’angle de déviation. Établir la relation
\begin{equation}
D = i + i' - A
\end{equation}
En appliquant le principe du retour inverse de la lumière, montrer que, pour une valeur de $D$ possible donnée, il existe deux couples de solutions ($i$, $i'$). En déduire l'égalité de $i$ et de $i'$ lorsque $D$ passe par un minimum (supposé unique).

\sol{
C'est également du cours. D'après le principe retour inverse de la lumière si $i$ en entrée et $i'$ en sortie donnent une déviation $D$, alors $i'$ en entrée conduira à $i$ en sortie et donc aussi une déviation $D$. Puisque le minimum est unique, il n'existe qu'un incidence pour laquelle  il est obtenu, donc nécessairement on a $i = i'$.
}

\item Déterminer la valeur $i_0$ de $i$ correspondant au minimum de déviation en fonction de $n$ et de $A$. \'Etablir une relation entre $n$, l'angle $A$ et la déviation minimale $D_m$.

\sol{
On a $\sin i_0 = n \sin r_0$ et d'autre part $A = 2 r_0$, donc on a :
\begin{equation*}
\sin i_0 = n \sin \left(\frac{A}{2}\right)
\end{equation*}
soit :
\begin{equation*}
i_0 = \mathrm{arcsin}\left(n \sin \left(\frac{A}{2}\right)\right)
\end{equation*}
Or on a $D_m = 2i_0 - A$, donc :
\begin{equation*}
\sin\left(\frac{D_m+A}{2}\right) = n \sin \left(\frac{A}{2}\right)
\end{equation*}
}

\item En déduire que l'indice $n$, les angles $A$ et $D_m$ vérifient une relation du type :
\begin{equation}
n = \frac{f\left(\frac{A+D_m}{2}\right)}{f\left(\frac{A}{2}\right)}
\end{equation}
où $f$ est une fonction que l'on précisera.

\sol{
D'après la question précédente :
\begin{equation*}
n = \frac{\sin\left(\frac{A+D_m}{2}\right)}{\sin\left(\frac{A}{2}\right)}
\end{equation*}
La fonction $f$ cherchée est donc la fonction sinus.
\item Si on trace $n$ en fonction de $1/\lambda^2$ alors les points s'alignent sur une droite et on trouve :
\begin{equation*}
a = 1.713 \qquad b = 0.01479\ \mu\mathrm{m}^2
\end{equation*}
Le milieu est dispersif car l'indice dépend de la longueur d'onde.
}

% \bonus{
% \item On éclaire le prisme avec une lampe à vapeur de mercure, pour laquelle on a mesuré $D_m$ pour différentes longueurs d'onde $\lambda$ et obtenu les valeurs de $n$ correspondantes, reprise dans la table ci-dessous.

% \begin{center}
% \begin{tabular}{|c|c|c|c|c|c|}
% \hline
% $\lambda$ (\micro m) & 0.4047 & 0.4358 & 0.4916 & 0.5461 & 0.5770 \\
% \hline
% $n$ & 1.803 & 1.791 & 1.774 & 1.762 & 1.757 \\
% \hline
% $\lambda^{-2}$ (\micro m$^{-2}$) & 6.11 & 5.27 & 4.14 & 3.35 & 3.00 \\
% \hline
% \end{tabular}
% \end{center}

% Montrer que $n$ peut se mettre sous la forme
% \begin{equation*}
% n = a + \frac{b}{\lambda^2}
% \end{equation*}
% où $a$ et $b$ sont des constantes. Est-ce que le prisme est dispersif et pourquoi ?

% \sol{
% Si on trace $n$ en fonction de $1/\lambda^2$ alors les points s'alignent sur une droite et on trouve :
% \begin{equation*}
% a = 1.713 \qquad b = 0.01479\ \mu\mathrm{m}^2
% \end{equation*}
% Le milieu est dispersif car l'indice dépend de la longueur d'onde.
% }

% \item Pour une lampe à vapeur de cadmium, on mesure un indice égal à 1.777. En déduire la longueur d'onde et donner la couleur correspondante.

% \sol{
% \'A partir de $n$, la longueur d'onde est obtenue par :
% \begin{equation*}
% \lambda = \sqrt{\frac{b}{n-a}} = 480.7\ \mathrm{nm}
% \end{equation*}
% La raie du cadmium étudiée est donc bleue.
% }
% }
\end{quest}
\end{exo}
%%%%%%%%%%%%%%%%%%%%%%%%%%%%%%%%%%%%%%%%%%%%%%%%%%%%%

%%%%%%%%%%%%%%%%%%%%%%%%%%%%%%%%%%%%%%%%%%%%%%%%%%%%%%%%%%
\begin{exo}[PCSI]{Voir la muraille de Chine depuis la Lune}
Certaines rumeurs disent que que la muraille de Chine est le seul bâtiment construit par
l'homme qui est visible par les astronautes depuis la lune (400 000 km). Est-ce vrai ?
Données (utiles ?) : largeur de la muraille de Chine - environ 10 m; limite de résolution de l'oeil (angle en dessous duquel on ne peut distinguer 2 objets vu sous ce même angle) - 1,5' d'arc.

\sol{
Pour répondre à la question, il faut d'abord réfléchir d'un point de vue moins ... astronomique. Prenons
les graduations d'une règle. Les graduations millimétriques d'une règle deviennent invisibles à une distance de
l'ordre de 10m. Un théorème de Pythagore nous donne alors un angle de vue d'environ $\alpha=10^{-4}$ rad.
L'angle sous lequel on voit la largeur de la muraille depuis la lune fait $10/4.10^8=2.10^{-8}$ rad.
Ceci est bien en dessous de la limite de résolution de l'oeil. La rumeur est complètement fausse ..
}

\end{exo}

%%%%%%%%%%%%%%%%%%%%%%%%%%%%%%%%%%%
\begin{exo}[PCSI][3][colle]{Réfraction sur un bassin}
Un observateur de taille $t=1,8$ m regarde un objet au fond d’un bassin depuis le bord. Le
bassin de hauteur $h=1,5$ m contient une épaisseur $e=1$ m d’eau (indice $n=1,33$).
L’objet semble situé à une distance $d=1$ m du bord. Déterminer la véritable position de
l’objet.
\begin{center}
  \includegraphics[scale=0.5]{vision_objet_eau.JPG}
 \end{center}

\sol{
L’observateur a l’impression que l’objet est dans la direction du rayon qu’il reçoit. On voit sur la
figure ce-dessus que le rayon provenant de l’objet s’est éloigné de la normale au dioptre (air-eau) puisque
$n > 1$. On note $A$ la position réelle de l’objet et $B$ sa position apparente.
\begin{itemize}
    \item Dans le triangle $TOB$ : $d=OB=(h+t)\mathrm{tan}(i)$ d'où $i=16,9^{\circ}$.
    \item Dans le triangle $TO^{'}I$ : $O^{'}I=(h-e+t)\mathrm{tan}(i)$
    \item Dans le triangle $II^{'}A$ : $AI^{'}=e \ \mathrm{tan}(r)$
    \item Loi de Descartes en $I$ : $\mathrm{sin}(i)=n\ \mathrm{sin}(r)$
\end{itemize}
D'autre part, les angles étant inférieurs à $\pi/2$, on a :
\begin{equation*}
\mathrm{cos}(r)=\sqrt{1-\mathrm{sin}^2(r)}=\sqrt{1-\left(\frac{\mathrm{sin}(i)}{n}\right)^2}
\end{equation*}
Donc :
\begin{equation*}
    OA=O^{'}I + I^{'}A=e \ \mathrm{tan}(r)+(h-e+t)\mathrm{tan}(i)
\end{equation*}
On a alors :
\begin{equation*}
    OA=e\frac{\mathrm{sin}(i)}{n}\frac{1}{\sqrt{1-\left(\frac{\mathrm{sin}(i)}{n}\right)^2}}+(h-e+t)\mathrm{tan}(i)=0,92 \mathrm{m}
\end{equation*}
}

\end{exo}

%%%%%%%%%%%%%%%%%%%%%%%%%%%%%%%%%%%%%%%%%%%%%%%%%%%%%
\begin{exo}[PCSI]{Fibre optique}
Une fibre optique est formée d'une âme en verre d'indice optique $n_{c} = 1.66$ entourée d'une gaine d'indice optique $n_{g} = 1.52$.

\begin{quest}
 \item Quelle est la valeur maximale de l'angle $i$ d'incidence pour que la lumière soit transmise sans perte le long de la fibre ?
 \item Justifier que si $i_{lim} \leq i \leq i_{lim}$ la propagation se fait sans perte d'énergie.
\end{quest}
\end{exo}
%%%%%%%%%%%%%%%%%%%%%%%%%%%%%%%%%%%%%%%%%%%%%%%%




%%%%%%%%%%%%%%%%%%%%%%%%%%%%%%%%%%%%%%%%%%%%%%%%%%%%%
\begin{exo}[PCSI]{Image par une lentille convergente - 2}
 On objet $\rm AB$ de 1 cm est placé à 5 cm en amont d'une lentille convergente de focale 2 cm et de centre optique $\rm O$.

 \begin{quest}
  \item Calculer la vergence de la lentille.
  \item Construire l'image $\rm A'B'$ dans le cas où $\rm AB$ est perpendiculaire à l'axe optique. Qu'en est-il si l'objet est incliné ?
  \item Calculer les distances $A'B'$ et $OA'$.
\item Définir le grandissement transversal $\gamma$. Le calculer.
 \end{quest}
\end{exo}
%%%%%%%%%%%%%%%%%%%%%%%%%%%%%%%%%%%%%%%%%%%%%%%%%%%%%%%%%

%%%%%%%%%%%%%%%%%%%%%%%%%%%%%%%%%%%%%%%%%%%%%%%%%%%%%%%%%%%
\begin{exo}[PCSI]{Incidence de Brewster}
Considérons un dioptre entre de l’air (d’indice de réfraction égal à $1$) et un milieu d'indice de réfraction $n$. Il existe un angle d'incidence $i_{B}$, dit de Brewster, tel que l'angle
entre les rayons réfléchi et réfracté émergent dans des directions perpendiculaires.

\begin{quest}
 \item Faire un schéma de la situation.
 % lors d'un rayon incident à l'angle de Brewster. Tracer également les rayons

 \sol{
\begin{center}
 \includegraphics[scale=1]{brewster_angle}
 \end{center}
 }
 \item Rappeler les lois de Snell-Descartes de la réflexion et de la réfraction.
 \item Calculer $i_{B}$. Application numérique pour l’eau ($n=\num{1.33}$) et le verre ($n=\num{1.5}$).
 \sol{
 On a $\pi/2 = \theta_{\rm rflx} + \theta_{\rm rfrc}$ or   $\theta_{\rm rflx} =i$ et $n \sin \theta_{\rm rfrc} =  \sin i = n \cos i$ et donc  $\tan i = n$.
 }
\end{quest}
\end{exo}
%%%%%%%%%%%%%%%%%%%%%%%%%%%%%%%%%%%%%%%%%%%%%%%%%%%%%%%%%%%%

%%%%%%%%%%%%%%%%%%%%%%%%%%%%%%%%%%%%%%%%%%%%%%%%%%%%%%%%%%%
\begin{exo}[PCSI]{Indice d'un milieu}
Considérons un cube de matière transparent environné d'air. Un rayon lumineux a un angle d'incidence de $45^{\circ}$ sur une des faces du cube. Il émerge sur un coté
perpendiculaire au premier puis ressort dans l'air en rasant la face de sortie.

\begin{quest}
 \item Faire un schéma.
 \item Calculer l'indice $n$ du milieu transparent.
\end{quest}
\end{exo}
%%%%%%%%%%%%%%%%%%%%%%%%%%%%%%%%%%%%%%%%%%%%%%%%%%%%%%%%%%


%%%%%%%%%%%%%%%%%%%%%%%%%%%%%%%%%%%%%%%%%%%%%%%%%%%%%%%%%%%%%
\begin{exo}[PCSI][2]{Lunette astronomique}
	\Notions{\item Optique géométrique \task Conditions de Gauss \task Doublet de lentille \task Système afocal}
Une lunette astronomique est constituée d'un objectif $\mc{L}_{1}$  de distance focale $f'_{1}=\SI{50}{cm}$ et d'un oculaire $\mc{L}_{2}$ de distance focale $f'_{2}=\SI{2}{cm}$. On note $d=\overline{O_1O_2}$ la distance entre les centres de $\mc{L}_1$ et $\mc{L}_2$.

\begin{quest}
 \item Comment régler la distance entre les lentilles pour que le système soit afocal ? Schématiser le doublet de lentilles et représenter la marche de deux rayons particuliers. Quelle est la différence entre  une lunette astronomique et la lunette de Galilée ?
 \sol{
	\begin{enumerate}
		\item 
		Considérons un point $A$ à l'infini sur l'axe optique. L'image par l'objectif de $A$, noté $A_1$, se trouve donc en $F_1'$. Or $A_1$ doit également se situer en $F_2$ pour que son image $A'$ par l'oculaire soit rejetée à l'infini, de sorte que l'œil puisse observer un objet lointain sans accommoder. Ainsi $F_1'=F_2$. On peut représenter le système optique par le schéma synoptique suivant,

		\eqbox{
			\infty  \underset{\mc{L}_1}{\rightarrow} F_1' = F_2 \underset{\mc{L}_2}\rightarrow \infty .
			}
		Ce système ne possède pas de foyer image (l'image d'un objet à l'infini est elle-même à l'infini) ni objet (l'objet donnant une image à l'infini est lui-même à l'infini). On parle donc de \fbox{système afocal}. La distance entre les lentilles vaut donc 
			\eqbox{
				d=f_1'+f_2' = 52 \ \mathrm{cm}.
				}  
	\item  On se place dans les \udl{conditions de Gauss} : les rayons sont paraxiaux, i.e. peu inclinés par rapport à l'axe optique ($\theta\ll1$) ; et proches de l'axe optique (le rayon doit frapper la lentille près du centre). Pour tracer les rayons, on considère un point objet $B$ à l'infini hors de l'axe optique. On rappelle les trois rayons particuliers que l'on peut utiliser pour tracer la marche des rayons en sortie d'une lentille :
	\begin{itemize}
		\item celui qui passe par le foyer objet $F$ sort parallèlement à l'axe optique ;
		\item celui qui passe par le centre optique $O$ sort sans être dévié ;
		\item celui qui arrive parallèlement à l'axe optique passe par le foyer image $F'$.
	\end{itemize}
	\begin{center}
\includegraphics[scale=1]{lunette_astro}
\end{center}
\item
La lunette de Galilée est aussi constituée d'un doublet de lentilles minces, mais elle n'est \fbox{pas afocale} et permet d'observer des objets proches (terrestres). Elle est constituée d'un objectif convergent $\mc{L}_1$ et d'un oculaire divergent $\mc{L}_2$.
\end{enumerate}
}
 \item Définir le grossissement angulaire $G$. Quelle est son expression en fonction des distances focales $f'_{1}$ et $f'_{2}$ ? Le
  calculer numériquement dans ces conditions. Le grandissement transversal $\gamma$ est-il défini ?
  \sol{
	\begin{enumerate}
		\item Le \udl{grossissement angulaire} $G$ a pour définition 
		\eqbox{
		G = \frac{\theta'}{\theta},
		}
		où $\theta'$ est l'angle formé par le rayon sortant de l'oculaire et l'axe optique, et $\theta$ est l'angle formé par le rayon incident sur l'objectif et l'axe optique, pour un faisceau incident parallèle. 
		\item Pour déterminer $G$ en fonction des distances focales, on utilise des relations de trigonométrie : 
		$$
		\tan \theta = \dfrac{\overline{F_2 B_1}}{\overline{O_1F_1'}} = \dfrac{\overline{F_2 B_1}}{f_1'}  \text{ et } \tan \theta' = \dfrac{\overline{F_2 B_1}}{\overline{O_2 F_2'}}  = -\dfrac{\overline{F_2 B_1}}{f_2'}.
		$$
		\udl{Remarque} : les distances $\overline{F_2 B_1}$, $\overline{O_1F_1'}$, et $\overline{O_2 F_2'}$ sont algébriques : elles peuvent être positives ou négatives, de même que les angles $\theta$ et $\theta'$. Penser à contrôler systématiquement le signe des relations ! Sur la figure, on a $\theta >0$ et $\theta' <0$.


		Dans les \udl{conditions de Gauss}, $\tan \theta \simeq \theta$ et $\tan \theta' \simeq \theta'$,
		et donc,
		$$
		G=\frac{\theta'}{\theta} \simeq \frac{\tan \theta'}{\tan \theta} \simeq -\dfrac{\overline{F_2 B_1}}{f_2'} \times \frac{f_1'}{\overline{F_2 B_1}},
		$$
		soit 
		\eqbox{
		G \simeq -\frac{f_1'}{f_2'}.
		}
		\AN~: \fbox{$G= -25$}. Le signe négatif indique que l'image est renversée.

		\item Le \udl{grandissement transversal} n'est pas défini pour un système afocal, puisque l'objet $AB$ et l'image $A'B'$ sont de taille infinie.
	\end{enumerate}
		
}
 \item La lunette donne une image renversée. Pour redresser l'image, on intercale entre $\mc{L}_{1}$ et $\mc{L}_{2}$ une lentille $\mc{L}_{3}$ de distance focale $f'_{3}=\SI{2}{cm}$. Le système des trois lentilles est toujours afocal.
  \begin{enumerate}
   \item Quelle est la position de $\mc{L}_{3}$ ?
   \sol{
	   Notons $A_1$ l'image de $A$ par $\mc{L}_1$ et $A_3$ l'image de $A_1$ par $\mc{L}_3$. Le système est afocal, donc $\infty \rightarrow \infty$ par le système constitué du triplet de lentille. On a donc le schéma synoptique suivant,
	%    $$
	\eqbox{
		\infty \underset{\mc{L}_1}{\rightarrow} F_1' \underset{\mc{L}_3}\rightarrow F_2 \underset{\mc{L}_2}\rightarrow \infty.
		}
	%    $$
	   Le foyer $F_1'$ est donc le conjugué du foyer $F_2$ par $\mc{L}_3$. Pour trouver la position relative de $F_2$ par rapport à $F_1'$, on utilise la relation de \udl{conjugaison de Descartes} pour la lentille $\mc{L}_3$,
	   $$
	   \frac{1}{f_3'} = \frac{1}{\overline{O_3F_2}} - \frac{1}{\overline{O_3F_1'}}.
	   $$
	   Il existe une \fbox{infinité de positions} de $\mc{L}_3$ qui convienne si l'on autorise à bouger $\mc{L}_2$ également, ce qui est nécessaire puisque $F_1'$ et $F_2$ ne sont plus confondus. 

	   Pour réduire l'encombrement au maximum, on peut choisir la position de $\mc{L}_3$ qui minimise la distance objet-image, i.e. la distance $\overline{F_1'F_2}$. On peut montrer qu'il faut placer pour cela la lentille $\mc{L}_3$ à mi-distance entre $F_1'$ et $F_2$, auquel cas 
	   $$
	    \overline{O_3F_2} = \overline{F_1'O_3} = 2f_3'.
	   $$
	   On trouve alors la distance entre l'objectif et la lentille $\mc{L}_3$,
	   \eqbox{
	   \overline{O_1O_3} = \overline{O_1F_1'} + \overline{F_1'O_3} = f_1' + 2f_3' = \SI{54}{cm}.
	   }
	   \udl{Remarque} : la longueur complète de la lunette vaut alors
	   $$
	   	\overline{O_1O_2} = \overline{O_1O_3} + \overline{O_3O_2} = \overline{O_1O_3} + \overline{O_3F_2} + \overline{F_2O_2} = f_1' + 4f_3' + f_2' = \SI{60}{cm}.
	   $$
	   }
%    \item Quel vaut l'encombrement  de la lunette ?
   \item Quel est le nouveau grossissement $G$ ?
   \sol{
	Il faut faire un schéma de la situation et utiliser des relations de trigonométrie.On schématise le système optique ci-dessous dans la situation discuté dans la question précédente, où l'encombrement est minimal.
	\begin{fig}[label]
		\centering
		\includegraphics[scale=0.9]{triplet_astro}
		%\capt{}
	\end{fig}
	On remarque que l'image n'est plus renversée, et donc le grossissement $G$ est positif.
   }
  \end{enumerate}
\end{quest}
\end{exo}
%%%%%%%%%%%%%%%%%%%%%%%%%%%%%%%%%%%%%%%%%%%%%%%%%%%%%%%%%%%%%%%%%%


%%%%%%%%%%%%%%%%%%%%%%%%%%%%%%%%%%%%%%%%%%%%%%%%%%%%%%%%%%%%%%%%%%%%
\begin{exo}[PCSI][1][colle]{Étude d'un spectroscope à prisme}
Un rétroprojecteur est constitué d'une lentille convergente de focale f', d'un miroir faisant un angle de $45^{\circ}$ avec l'axe optique et d'un écran. On éclaire uniformément
un onjet $AB$. On note $d=10cm$ la distance lentille-miroir, $D$ la distance miroir-écran et $h$ la distance objet-lentille. $h$ peut varier entre 30 et 42 cm.

\begin{quest}
 \item Donner un schéma rectiligne équivalent du système.
 \item Dans quel sens doit-on placer le document pour qu'il apparaisse dans le bon sens à l'écran ?
 \item La mise au point est faite pour $h=33$ cm et $D=2.3$ m.
 \begin{quest}
  \item Quelle est la focale de la lentille ?
  \item Quel est le grandissement ?
  \item Le document ayant une largeur de 21 cm, quelle est sa largeur à l'écran ?
 \end{quest}
 \item On garde la même lentille et le même document. L'écran est large de 2.1 m. On considère que le montage est optimal si l'image occupe tout l'écran. Calculer $D$ puis $h$.
\end{quest}
\end{exo}
%%%%%%%%%%%%%%%%%%%%%%%%%%%%%%%%%%%%%%%%%%%%%%%%%%%%%%%%%%%%%%%%%%%%%%%%




%%%%%%%%%%%%%%%%%%%%%%%%%%%%%%%%%%%%%%%%%%%%%%%%%%%%%%%%
\begin{exo}[PCSI]{Incidence de Brewster et polarisation}

Un rayon lumineux arrive à l'interface plane séparant l'air d'un milieu d'indice $n$. Il se scinde en un rayon réfléchi et un rayon réfracté.
\begin{quest}
\item Trouver l'angle d'incidence $i_B$, appelé angle de Brewster, pour lequel ces deux rayons sont perpendiculaires entre eux. Faire l'application numérique dans le cas de l'eau d'indice 1.33 puis du verre d'indice 1.5.
\sol{
Si on appelle $i$ l'angle d'incidence et $r$ l'angle de réfraction alors que pour que les deux rayons soient perpendiculaires il faut que :
\begin{equation*}
r = \frac{\pi}{2} -i
\end{equation*}
Or on a $\sin i = n \sin r$ soit $\sin i = n \cos i$ et donc $\tan i = n$
L'angle d'incidence de Brewster est :
\begin{equation*}
\boxed{i_B = \mathrm{arctan}(n)}
\end{equation*}
L'application numérique donne $53^o$ pour l'eau et $56^o$ pour le verre.
}
\item Lorsque l'angle d'incidence est $i_B$, la lumière réfléchie est polarisée rectilignement selon la direction perpendiculaire au plan d'incidence. Quelle application du polariseur pouvez-vous imaginer à partir de cette propriété ?
\sol{
En munissant l'objectif d'un polariseur, on peut réduire voire éliminer (si l'incidence est d'incidence de Brewster) les reflets sur l'eau ou sur une vitre en orientant convenablement le polariseur.
}
\end{quest}
\end{exo}
%%%%%%%%%%%%%%%%%%%%%%%%%%%%%%%%%%%%%%%%%%%%%%%%%%%%%%%%%%

%%%%%%%%%%%%%%%%%%%%%%%%%%%%%%%%%%%%%%%%%%%%%%%%%%%%%%%%%%%%%
\begin{exo}[PCSI]{Dispersion de la lumière solaire lors d'un arc-en-ciel}
Un milieu dispersif est un milieu dont l'indice optique dépend de la longueur d'onde de la lumière. Dans ce problème on s'intéresse à une conséquence du caractère dispersif de l'eau : le phénomène d'arc en ciel. L'indice optique de l'eau est donnée par la loi de Cauchy :
\begin{equation*}
n = a + \frac{b}{\lambda^2}
\end{equation*}
où $a$ et $b$ sont des constantes positives caractéristiques du milieu et $\lambda$ la longueur d'onde.
\begin{quest}
\item \textit{Questions préliminaires}
\begin{quest}
\item Rappeler les lois de Descartes pour la réfraction d'un rayon lumineux passant de l'air (milieu d'indice unité) vers un milieu d'indice $n$. On fera un schéma en notant $i$ l'angle d'incidence et $r$ l'angle de réfraction. Exprimer la dérivée $\frac{\mathrm{d}r}{\mathrm{d}i}$ exclusivement en fonction de l'indice $n$ et du sinus de $i$.

\sol{
Loi de Descartes pour la réfraction : le rayon refracté appartient au plan d'incidence et
\begin{equation*}
\sin i = n \sin r
\end{equation*}
On dérive cette loi par rapport à $i$ sachant que $n$ est indépendant de $i$ :
\begin{equation*}
\cos i = n \cos r \frac{\mathrm{d}r}{\mathrm{d}i} \qquad \Rightarrow \qquad \frac{\mathrm{d}r}{\mathrm{d}i} = \frac{1}{n}\frac{\cos i}{\cos r}
\end{equation*}
Les angles sont dans l'intervalle $\left[-\frac{\pi}{2}\ ;\ \frac{\pi}{2}\right]$, donc le cosinus est positif, on peut alors utiliser la formule $\cos x = \sqrt{1 - \sin^2 x}$ pour trouver :
\begin{equation*}
\frac{\mathrm{d}r}{\mathrm{d}i} = \frac{1}{n}\frac{\sqrt{1-\sin^2 i}}{\sqrt{1-\sin^2 r}} =  \frac{1}{n}\frac{\sqrt{1-\sin^2 i}}{\sqrt{1-\left(\frac{\sin i}{n}\right)^2}}
\end{equation*}
}
\item Exprimer, en fonction de $i$ et de $r$, la valeur de la dérivation du rayon lumineux, définie par l'angle entre la direction incidente et la direction émergente, orientées dans le sens de propagation.

\sol{
La déviation est l'angle dont a tourné le rayon lumineux par rapport à la direction incidente. Dans le cas de la réfraction, la déviation $D$ est égale à $r-i$.
}

\item Exprimer aussi, à l'appui d'un schéma, la déviation d'un rayon lumineux dans le cas d'une réflexion.

\sol{
Dans le cas de la réflexion la déviation vaut $D = \pi - 2i$.
}
\end{quest}
\item Lorsque le soleil illumine un rideau de pluie, on peut admettre que chaque goutte d'eau se comporte comme une sphère réceptionnant un faisceau de rayons parallèles entre eux.

Dans tout ce qui suit, on considérera que l'observation est faite par un {\oe}il accommodant à l'infini, c'est-à-dire assimilable à une lentille convergente (cristallin) capable de focaliser sur un écran (rétine) tout faisceau de lumière parallèle issu d'une goutte d'eau.

On cherche dans un premier temps les conditions pour que la lumière émergente, issue d'une goutte d'eau, se présente sous forme d'un faisceau de lumière parallèle. Pour cela on fait intervenir l'angle de déviation $D$ de la lumière à travers la goutte d'eau, mesuré entre le rayon émergeant et le rayon incident. Cet angle de déviation $D$ est une fonction de l'angle d'incidence $i$.

Exprimer la condition de parallélisme des rayons émergents en la traduisant mathématiquement au moyen de la dérivée $\frac{\mathrm{d}D}{\mathrm{d}i}$.

\sol{
Les rayons émergents sont parallèles quel que soit l'angle d'incidence $i$ si la déviation $D$ est indépendante de $i$ soit :
\begin{equation*}
\frac{\mathrm{d}D}{\mathrm{d}i} = 0
\end{equation*}
}

\bigskip

Une goutte d'eau quelconque, représentée par une sphère de centre $O$ est atteinte sous des incidences $i$ variables comprises entre 0 et 90$^o$. Son indice, pour une radiation donnée, sera noté $n$ tandis que celui de l'air sera pris égal à 1.
On considère les trois cas ci-dessous :
\begin{itemize}
\item lumière directement transmise,

\sol{
 LE triangle $OIJ$ est isocèle, dans $\alpha = -r$. Dans ce cas, la loi de Descartes en $J$ s'écrit $\sin\beta = n \sin (-r) = - n \sin r = - \sin i$. Donc :
\begin{equation*}
\beta = -i
\end{equation*}
}
\item lumière transmise après une réflexion partielle à l'intérieur de la goutte,

\sol{
La déviation en $I$ est $ D = r-i$, et celle en $J$ est $D' = \beta - \alpha = r-i$, soit globalement :
\begin{equation*}
D_1 = D + D' = 2 (r-i)
\end{equation*}
}
\item lumière transmise après deux réflexions partielle à l'intérieur de la goutte.

\sol{
On calcule la dérivée de $D_1$ par rapport à $i$ :
\begin{equation*}
\frac{\mathrm{d}D_1}{\mathrm{d}i} = 2 \frac{\mathrm{d}r}{\mathrm{d}i} - 2 = 2 \frac{1}{n}\frac{\sqrt{1-\sin^2 i}}{\sqrt{1-\left(\frac{\sin i}{n}\right)^2}} - 2
\end{equation*}
La dérivée est nulle si $\sqrt{1-\sin^2 i} = \sqrt{1-\left(\frac{\sin i}{n}\right)^2}$. En élevant au carré et en simplifiant on trouve que $n=1$ ce qui n'a pas d'intérêt. Ainsi, il n'est pas possible d'avoir des rayons parallèles émergents.
}

\end{itemize}
Les deux premiers sont représentés ci-dessous, on dessinera le troisième avec pour angles supplémentaires de réflexion $\phi$ et $\xi$ :

\begin{center}
\includegraphics[width = 0.7\textwidth]{Gouttes.jpg}
\end{center}

\item Pour le cas 1 :
\begin{quest}
\item Exprimer en fonction de l'angle d'incidence $i$ ou de l'angle de réfraction $r$, tous les angles marqués de lettres grecques.

\sol{
 L'application de la loi de Descartes pour la reflexion en $J$ donne $\beta = - \alpha$. D'autre part les triangles $OIJ$ et $OJK$ sont isocèles sont $ R = \beta = - \alpha = - \gamma$. Alors l'application de la loi de Descartes en $K$ donne $\sin \delta = n \sin \gamma = - n \sin r = - \sin i$, soit $\delta = -i$.
}

\item En déduire l'angle de déviation $D$ en fonction de $i$ et de $r$.

\sol{
Pour la déviation, on a la somme de trois termes : ladéviation en $I$, $D = r-i$ : la déviation en $J$, $D'= \pi - 2 \alpha$ ; la déviation en $K$, $D'' = \delta - \gamma = -i + r$. Donc :
\begin{equation*}
D_2 = D + D' + D'' = \pi + 4r - 2i
\end{equation*}
}
\item rechercher ensuite, si elle existe, une condition d'émergence d'un faisceau parallèle, exprimée par une relation entre $\sin i$ et l'indice $n$ de l'eau.

\sol{
On dérive $D_2$ :
\begin{equation*}
\frac{\mathrm{d}D_2}{\mathrm{d}i} = 4 \frac{\mathrm{d}r}{\mathrm{d}i} - 2 = 4 \frac{1}{n}\frac{\sqrt{1-\sin^2 i}}{\sqrt{1-\left(\frac{\sin i}{n}\right)^2}} - 2
\end{equation*}
Cette dérivée s'annule pour :
\begin{equation*}
\sin^2 i = \frac{4 - n^2}{3}
\end{equation*}
}

\end{quest}
\item Mêmes questions pour le cas 2.
\sol{
De la même manière : $r = \beta = \delta = -\alpha = - \gamma = - \phi$ et $\xi = -i$.
}


\item Mêmes questions pour le cas 3.


\sol{
On trouve alors :
\begin{equation*}
D_3 = 2\pi + 6r - 2 i
\end{equation*}
}

\item Le soleil étant supposé très bas sur l'horizon,, normal au dos d'un observateur, montrer que celui-ci ne pourra observer la lumière transmise que si la goutte d'eau se trouve sur deux cônes d'axes confondus avec la direction solaire et de demi-angles au sommet $\theta_2 = 180^o + D_2 = 180^o - |D_2|$ (justification de l'arc primaire) et $\theta_2 = D_3 -180^o$ (justification de l'arc secondaire).


\sol{
Et donc en dérivant :
\begin{equation*}
\frac{\mathrm{d}D_3}{\mathrm{d}i} = 6 \frac{\mathrm{d}r}{\mathrm{d}i} - 2 = 6 \frac{1}{n}\frac{\sqrt{1-\sin^2 i}}{\sqrt{1-\left(\frac{\sin i}{n}\right)^2}} - 2
\end{equation*}
Cette dérivée s'annule pour :
\begin{equation*}
\sin^2 i = \frac{9 - n^2}{8}
\end{equation*}
}

\sol{
Le soleil étant très bas sur l'horizon, on peut supposer que les rayons incidents sont horizontaux. La figure ci-dessous représente les deux cas possible.
\begin{center}
\includegraphics[width = 0.6\textwidth]{Arc-en-ciel.jpg}
\end{center}
Les angles $\theta_2$ et $\theta_3$ entre les rayons incidents pour l'{\oe}il et les rayons provenant du soleil vérifient :
\begin{equation*}
\theta_2 = \pi + D_2 = \pi - |D_2|
\end{equation*}
\begin{equation*}
\theta_3 = D_3 - \pi
\end{equation*}
}

\item Les angles $\theta_2$ et $\theta_3$ dépendant de l'indice $n$ de l'eau, on observe un phénomène d'irisation dû au fait que cet indice évolue en fonction de la longueur d'onde.

Calculer ces angles pour le rouge et le violet, sachant que pour le rouge l'indice vaut 1.3317 tandis que pour le violet il est égal à 1.3448.

En admettant que l'observateur se trouve face à un rideau de pluie, dessiner la figure qui apparait dans son plan d'observation en notant la position respective des rouges et des violets.

\sol{
Toutes les relations pour les cas 2 (arc primaire) et 3 (arc secondaire) ont été vu dans les questions précédentes. On calcule $i$ avec la relation (11) ou (14) On trouve (les angles sont en radians sauf dans la dernière colonne) :

\begin{center}
\begin{tabular}{|c|c|c|c|c|c|}
\hline
Couleur & $i$ & $r$ & $|D_2|$ & $\theta_2$ & $\theta_2$ ($^o$) \\
\hline
rouge & 1.0382 & 0.7035 & 2.4039 & 0.7377 & 42.268 \\
\hline
violet & 1.0250 & 0.6887 & 2.4366 & 0.7050 & 40.393 \\
\hline
\hline
Couleur & $i$ & $r$ & $D_3$ & $\theta_3$ & $\theta_3$ ($^o$) \\
\hline
rouge & 1.2546 & 0.7948 & 4.0238 & 0.8823 & 50.549 \\
\hline
violet & 1.2473 & 0.7825 & 4.0830 & 0.9414 & 53.936 \\
\hline
\end{tabular}
\end{center}


L'angle $\theta_2$ est plus grand pour le rouge que pour le violet. Les rayons rouges observés par l'{\oe}il proviennent donc de gouttes de plus haute altitude, car dans ce cas les rayons violets passent au dessus de l'{\oe}il, et ce sont les  rayons violets des gouttes de plus basse altitude qui entrent dans l'{\oe}il. On voit donc le violet en dessous du rouge pour l'arc primaire.

En ce qui concerne l'arc secondaire, il n'est pas toujours visible car beaucoup moins intense (car deux réflexions). On remarque que $\theta_3 > \theta_2$, donc les gouttes donnant l'arc secondaire visible par l'{\oe}il seront à plus haute altitude que celles donnant l'arc primaire.On remarque également que les couleurs sont inversées.

}
\end{quest}
\end{exo}
%%%%%%%%%%%%%%%%%%%%%%%%%%%%%%%%%%%%%%%%%%%%%%%%%%%%%%

%%%%%%%%%%%%%%%%%%%%%%%%%%%%%%%%%%%%%%%%%%%%%%%%%%%%%%
\begin{exo}[PCSI]{Fibre optique à saut d'indice}%{1}
%Le guidage de la lumière peut être assuré par des fibres optiques.
\Notions{\item Optique géométrique \task Fibre optique \task Réflexion totale \task Loi de Snell-Descartes}
Une fibre optique est constituée d'un cylindre de verre appelé c{\oe}ur, de rayon $a$ et d'indice $n_c$, entouré d'une gaine transparente de rayon $b$ et d'indice de réfraction $n_g<n_c$.
La gaine contribue non seulement aux propriétés mécaniques de la fibre mais évite aussi les fuites de lumière vers d'autres fibres en cas de contact.
Le diamètre du c{\oe}ur d'une fibre varie de $\SI{3}{}$ à $\SI{200}{\mu m}$ selon ses propriétés, et le diamètre extérieur de la gaine peut atteindre $\SI{400}{\mu m}$.
L'ensemble, en particulier la face d'entrée, est au contact de l'air d'indice $n_0=1$.

\begin{fig}[label]
	\centering
	\includegraphics[width=0.7\columnwidth]{fibre_saut.png}
	%\capt{}
\end{fig}
\begin{quest}
\item Un rayon lumineux arrive en $O$, centre de la face d’entrée de la fibre. On appelle $i$ l'angle d'incidence. Déterminer en fonction de $n_c$ et $n_g$ la condition que doit satisfaire $i$ pour que le rayon réfracté se propage dans le c{\oe}ur et ne s'échappe pas par la gaine. Quelle condition doit vérifier l’angle $\alpha$ pour avoir réflexion totale ? On appelle angle d'acceptance $i_a$ de la fibre la valeur maximale de $i$. Déterminer l'expression de $i_a$.
\sol{
Le rayon est guidé le long de la fibre s'in y a reflexion totale sur la surface de séparation c{\oe}ur-gaine, ce qui est possible puisque $n_2>n_1$. Pour cela il faut que l'angle d'incidence par rapport à la normale, noté $\theta$, vérifie :
\begin{equation}
|\sin\theta| > \frac{n_2}{n_1}
\end{equation}
Or, si on note $r$ l'angle par rapport à l'axe du rayon réfracté en entrée de la fibre, on a :
\begin{equation}
\theta = r - \frac{\pi}{2}
\end{equation}
De plus on a :
\begin{equation}
n_1 \sin r = n_0 \sin i
\end{equation}
On a donc :
\begin{equation}
|\sin\theta| =\cos r = \sqrt{1-\cos^2r} = \sqrt{1-\left(\frac{n_0\sin i}{n_1}\right)^2}
\end{equation}
La condition de réflexion totale devient donc :
\begin{equation}
\frac{n_0^2}{n_1^2}\sin^2 i < 1 - \frac{n_2^2}{n_1^2}
\end{equation}
Soit $i_a < i < i_a$ avec :
\begin{equation}
\boxed{i_a = \mathrm{arcsin}\left(\frac{\sqrt{n_1^2 -n_2^2}}{n_0}\right)}
\end{equation}
}
\item On appelle ouverture numérique ON de la fibre la quantité ON$ = n_0 \sin i_a$. Exprimer ON en fonction de $n_c$ et $n_g$. Application numérique pour $n_c = \num{1.456}$ (silice) et $n_g = \num{1.410}$ (silicone).
\sol{
D'après la définition on a :
\begin{equation}
\boxed{ON = \sqrt{n_1^2 -n_2^2}}
\end{equation}
L'application numérique donne $ON = 0.363$.
}
\item Le signal transporté par la fibre est constitué d'impulsions lumineuses d'une durée $\tau$ à intervalle régulier $T$. Quelle valeur minimale de $T$ faut-il choisir pour que deux impulsions successives ne se chevauchent pas en sortie de la fibre ? Calculer la bande passante et la comparer avec celle d'un téléphone portable (64 Mbits par seconde) et celle de la télévision (100 Mbits par seconde).
\sol{

Le temps le plus court est celui obtenu en se déplaçant sur l'axe. On obtient comme valeur $\frac{Ln_1}{c}$. Le temps le plus long correspond à l'angle d'accptance $i_a$ soit $\frac{Ln_1}{c\cos\theta}$ On en déduit donc que :
\begin{equation}
\boxed{\delta \tau = \frac{Ln_1}{c}\left(\frac{1}{\cos\theta} -1\right)}
\end{equation}
L'application numérique donne $\delta\tau = 0.158$ $\mu$s.
Il faut que $T > \delta\tau$ soit $f < \frac{1}{\delta\tau} = \Delta f = 6.3\ 10^6$ bits par second. La fibre étudiée convient pour le téléphone mais pas pour la télévision.
}
\end{quest}
\end{exo}
%%%%%%%%%%%%%%%%%%%%%%%%%%%%%%%%%%%%%%%%%%%%%%%%


%%%%%%%%%%%%%%%%%%%%%%%%%%%%%%%%%%%
\begin{exo}[PC][1][colle]{Interférométrie à deux télescopes}
	En 1996, des astronomes ont déterminé, avec une excellente précision, la géométrie de l'étoile double Capella, dans le domaine spectral du proche infrarouge. Au lieu d'utiliser un seul télescope dont la pupille est percée de deux trous (Méthode de Fizeau), on couple deux télescopes identiques $\mathrm{T}_1$ et $\mathrm{T}_2$, de même ouverture circulaire, de diamètre négligeable par rapport à la ligne de base $a=\mathrm{T}_1 \mathrm{~T}_2$ (Fig. 2.5). On supposera la ligne de visée perpendiculaire à $\mathrm{T}_1 \mathrm{~T}_2$. On désigne par $D$ la distance $\mathrm{O} \Omega, x_1$ la position de l'étoile $\mathrm{E}_1$ selon un axe perpendiculaire à la direction $(O \Omega)$, et $x_2$ la position analogue de $\mathrm{E}_2$, avec $x_2=-x_1$. On se place dans le cas où $a=6,10 \mathrm{~m}$. En outre, le rayonnement est quasi-monochromatique et centré sur la longueur d'onde $\lambda=635 \mathrm{~nm}$. Enfin, l'étoile double est supposée symétrique : chaque étoile $\mathrm{E}_1$ et $\mathrm{E}_2$ produit la même intensité $I_0$.
	Données :
	- $(1+e)^n \approx 1+n e$ si $\epsilon \ll 1$;
	- $\cos p+\cos q=2 \cos \left(\frac{p+q}{2}\right) \cos \left(\frac{p-q}{2}\right)$.
	\begin{fig}[]
	  \centering
	  \includegraphics[width=0.4\textwidth]{2telescopes.png}
	  % \capt{Schéma de principe du dispositif.}
	\end{fig}
	\begin{quest}
	  \item Les deux étoiles $\mathrm{E}_1$ et $\mathrm{E}_2$ seront considérées comme deux sources ponctuelles. Ces deux sources sont-elles cohérentes? Justifier.
	  \item $\mathrm{E}_1$ envoie des rayons aux deux télescopes. Exprimer la différence de chemin optique
	  $$
	  \delta_1=\mathrm{E}_1 \mathrm{~T}_2-\mathrm{E}_1 \mathrm{~T}_1
	  $$
	  
	  On donnera le résultat en fonction de $\alpha, x_1$ et $D$ et l'on fera des approximations raisonnables de façon à simplifier le résultat.
	  - $15 / 85 \cdot$
	  2 Interference à deux ondes
	  \item On note $\delta_2$ la différence de chemin optique entre les rayons provenant de $\mathrm{E}_2$ et arrivant sur $\mathrm{T}_1$ et $\mathrm{T}_2$. En s'inspirant du calcul précédent, exprimer $\delta_2$.
	  \item Un dispositif annexe permet de faire interférer les ondes optiques issues des deux télescopes. Exprimer l'intensité résultant de l'interférence en fonction de $a, \lambda$ et $\theta$ l'angle sous lequel est vu l'étoile double.
	  \item Le dispositif interférentiel permet d'introduire une différence de marche supplémentaire $L$ variable. En déduire la nouvelle expression de l'intensité $I(L)$. Montrer que la fonction $I(L)$ varie périodiquement et que le contraste $\gamma=\frac{I_{\max }-I_{\min }}{I_{\max }+I_{\min }}$ est relie à $\theta$.
	\end{quest}
	\end{exo}
	%%%%%%%%%%%%%%%%%%%%%%%%%%%%%%%%%%%


%%%%%%%%%%%%%%%%%%%%%%%%%%%%%%%%%%%
\begin{exo}[PC][1][colle]{Déviation minimale des rayons dispersés par un réseau}
	On considère un réseau de fentes éclairé par une onde plane monochromatique (longueur d'onde $\lambda$ ) faisant un angle $\alpha$ avec la normale au réseau. On observe les rayons diffractés dans une direction repérée par l'angle de déviation $D$. Ces rayons interfèrent entre eux à l'infini.
	\begin{quest}
	  \item Exprimer la condition d'interférence constructive entre les rayons diffractés. En déduire la loi donnant la déviation correspondant aux pics de diffraction.
	  \item Montrer que pour un ordre $p$ fixé, lorsque l'on fait varier l'angle $\alpha$ la déviation passe par un minimum que l'on exprimera.
	  - $20 / 85$ $\cdot$
	  \item Quelle est la configuration du réseau par rapport aux faisceaux incident et diffracté d'ordre $p$, lorsque la déviation est minimale?
	\end{quest}
	%%%%%%%%%%%%%%%%%%%%%%%%%%%%%%%%%%%\subsection{Exercice 7.1}
Une ligne coaxiale est constituée de deux cylindres $C_{1}$ et $C_{2}$ infinis d'épaisseur très faible de rayons $a$ et $b$ parcourus par des courants d'intensités $I$ et $-I$. On note $V_{1}$ et $V_{2}$ les potentiels des deux cylindres et on néglige toute chute de tension le long de la ligne.\\
%%%%%%%%%%%%%%%%%%%%%%%%%%%%%%%%%%%%%%%%%%%
\begin{fig}[label]
\centering
\includegraphics[width=0.5\columnwidth]{2024_06_14_b4e4733f962f43b4e093g-1}
\capt{}
\end{fig}
%%%%%%%%%%%%%%%%%%%%%%%%%%%%%%%%%%%%%%%

\begin{quest}
  \item Déterminer l'expression de la puissance rayonnée dans la ligne coaxiale en fonction de $V_{1}, V_{2}$ et $I$. Conclure.

  \item Que dire de la puissance rayonnée à l'extérieur de la ligne coaxiale?

\end{quest}
\end{exo}
%%%%%%%%%%%%%%%%%%%%%%%%%%%%%%%%%%%%%%%%%%%%

